%% toc: P2 22 22.1,22.2
\section{The one-sorted reduction for signatures of every cardinality}
\label{sec:p02-w02b-p02-033-any-cardinality}

The reduction of \cref{thm:p02-many-sorted-translation} turns each sorted
variable into a one-sorted variable.  The one-sorted variable supply
\(\mathsf{Var}=\{x_n:n\in\N\}\) of \cref{def:p02-fol-signature} is countable,
while a first-order signature may have function and relation symbols of any
cardinality.  This section separates the two sources of size.  The number of
function and relation symbols plays no role in any step of the reduction.  The
number of sorts does: with uncountably many sorts there are uncountably many
sorted variables, and no renaming of them into \(\mathsf{Var}\) can be exact.
The section therefore gives a second form of the reduction, in which every
sorted variable that is free in a formula becomes a new constant symbol and
bound sorted variables are renamed afresh inside each formula.  That form
needs no injection of variables and holds for every strict signature.  The
variable form is then recovered from it whenever an injection of the sorted
variables into \(\mathsf{Var}\) exists, and the impossibility of the variable
form for uncountably many sorts is proved.

\emph{Inputs.}  A strict many-sorted signature
\(\Lang_{\mathrm{MS}}=(\mathsf{Sort},\mathcal F,\mathcal R,\operatorname{tp})\)
with the typing clauses, terms and well-sorted formulas of
\cref{def:p02-many-sorted-signature}, but with no bound on the cardinality of
\(\mathsf{Sort}\), \(\mathcal F\) or \(\mathcal R\); the strict structures and
the consequence relation \(\entails_{\mathrm{MS}}\) of
\cref{def:p02-many-sorted-structure}; the term values and satisfaction clauses
of \cref{def:p02-w02-p02-038-satisfaction}; the restriction
\(\mathcal A^{\downarrow}\) of \cref{def:p02-w02-p02-038-restriction}; and
one-sorted satisfaction and consequence \(\entails_{\mathrm{FOL}}\) over
nonempty domains (\cref{sec:fol-semantics-02}).  \emph{Conclusions.}  A
reduction of strict many-sorted consequence to one-sorted consequence with the
sort predicates \(Q_{\mathsf s}\), the sort-axiom families
\((\mathrm{Ne})\), \((\mathrm{Ct})\), \((\mathrm{Cl})\), \((\mathrm{Dj})\) of
\cref{def:p02-w02-p02-038-axiom-families}, no coverage axiom, and a coverage
variant for finitely many sorts, valid for every strict signature; the
satisfaction correspondence behind it in both directions; the variable form of
\cref{thm:p02-many-sorted-translation} and of
\cref{thm:p02-w02-p02-038-strict-correspondence,thm:p02-w02-p02-038-coverage-reduction}
for function and relation symbols of every cardinality whenever the sorted
variables inject into \(\mathsf{Var}\); and a proof that the variable form is
impossible when there are uncountably many sorts.  \emph{Ambient
assumptions.}  The metatheory is ZFC.  Foundation is used to name fresh
symbols and a fresh junk point.  No statement of this section uses the Axiom
of Choice.

\subsection{Sorted variables and parameter names}
\label{sec:p02-w02b-p02-033-parameters}

\begin{hbconvention}{def:p02-w02b-p02-033-variables}{Sorted variable stocks without a size bound}
Every sort \(\mathsf s\in\mathsf{Sort}\) has a countably infinite stock
\(V_{\mathsf s}\) of variables, the stocks of distinct sorts are disjoint, and
\(x^{\mathsf s}\in V_{\mathsf s}\) is the designated variable of
\(\mathsf s\) (\cref{sec:p02-w02-p02-038-clauses}).  Put
\(\mathrm{Var}_{\mathrm{MS}}\defeq\bigcup_{\mathsf s\in\mathsf{Sort}}V_{\mathsf s}\)
and write \(\operatorname{srt}(v)\) for the unique sort with
\(v\in V_{\operatorname{srt}(v)}\).  The sets \(\mathsf{Sort}\),
\(\mathcal F\) and \(\mathcal R\) may have any cardinality.  The untyped
symbols of \(\mathcal F\cup\mathcal R\), the predicates \(Q_{\mathsf s}\) of
\cref{def:p02-many-sorted-translation}, the set \(\mathsf{Var}\) and the
logical symbols are pairwise disjoint, as \cref{def:p02-fol-signature}
requires of the flattened signature \(\Lang_{\mathrm{MS}}^{\flat}\).
\index{many-sorted first-order logic!variable stock}
\end{hbconvention}

The map \(\mathsf s\mapsto x^{\mathsf s}\) is an injection of
\(\mathsf{Sort}\) into \(\mathrm{Var}_{\mathrm{MS}}\), because the stocks are
disjoint.  So \(\mathrm{Var}_{\mathrm{MS}}\) is uncountable whenever
\(\mathsf{Sort}\) is.

\begin{hbdefinition}{def:p02-w02b-p02-033-parameter-signature}{Parameter names and the parameter signature}
Let \(W\) be the union of \(\mathcal F\), \(\mathcal R\),
\(\{Q_{\mathsf s}:\mathsf s\in\mathsf{Sort}\}\), \(\mathsf{Var}\) and the set
of logical symbols of \cref{def:p02-fol-signature}.  For
\(v\in\mathrm{Var}_{\mathrm{MS}}\) the \emph{parameter name} of \(v\) is the
ordered pair \(\mathsf n_v\defeq(W,v)\).  The \emph{parameter signature}
\(\Lang_{\mathrm{MS}}^{\mathsf n}\) is the one-sorted signature obtained from
\(\Lang_{\mathrm{MS}}^{\flat}\) by adding every \(\mathsf n_v\) as a
constant symbol.  Distinct variables have distinct names, and no name lies in
\(W\) (\cref{thm:p02-w02b-p02-033-sentences}); so
\(\Lang_{\mathrm{MS}}^{\mathsf n}\) is a first-order signature in the sense of
\cref{def:p02-fol-signature}.
The \emph{parameter guard} is the set of one-sorted sentences
\[
 \mathsf{ParGuard}(\Lang_{\mathrm{MS}})\defeq
 \{Q_{\operatorname{srt}(v)}(\mathsf n_v):v\in\mathrm{Var}_{\mathrm{MS}}\}.
\]
\index{many-sorted first-order logic!parameter name}
\end{hbdefinition}

The parameter guard plays the part of
\(\mathsf{VarGuard}_{\mathrm{MS}}\) (\cref{def:p02-ms-variable-guard}): it
says that the object a sorted variable names lies in the right sort.  It is a
set of sentences, so it constrains structures, not assignments.

\begin{hbdefinition}{def:p02-w02b-p02-033-environment-translation}{Environments and the parameter translation}
An \emph{environment} is a finite injective partial function
\(\eta\) from \(\mathrm{Var}_{\mathrm{MS}}\) to \(\mathsf{Var}\); write
\(\dom\eta\) and \(\operatorname{ran}\eta\) for its domain and range.  For
\(v\in\mathrm{Var}_{\mathrm{MS}}\) let \(k(\eta)\) be the least \(k\in\N\)
with \(x_k\notin\operatorname{ran}\eta\), and let \(\eta\langle v\rangle\) be
the environment that sends \(v\) to \(x_{k(\eta)}\) and agrees with \(\eta\)
on \(\dom\eta\setminus\{v\}\).  For a term \(t\) of
\(\Lang_{\mathrm{MS}}\), the one-sorted term \(t^{\eta}\) replaces each
variable \(v\) of \(t\) by \(\eta(v)\) when \(v\in\dom\eta\) and by
\(\mathsf n_v\) otherwise, and leaves constants and function symbols in
place.  The \emph{parameter translation}
\(\operatorname{tr}^{\mathsf n}_{\eta}\) is defined by recursion on
well-sorted formulas:
\[
\begin{aligned}
 \operatorname{tr}^{\mathsf n}_{\eta}(P(t_1,\ldots,t_n))
   &\defeq P(t_1^{\eta},\ldots,t_n^{\eta}),\\
 \operatorname{tr}^{\mathsf n}_{\eta}(t=u)
   &\defeq(t^{\eta}=u^{\eta}),\\
 \operatorname{tr}^{\mathsf n}_{\eta}(\bottom)
   &\defeq\bottom,\\
 \operatorname{tr}^{\mathsf n}_{\eta}(\neg\varphi)
   &\defeq\neg\operatorname{tr}^{\mathsf n}_{\eta}(\varphi),\\
 \operatorname{tr}^{\mathsf n}_{\eta}(\varphi\mathbin{\square}\psi)
   &\defeq\operatorname{tr}^{\mathsf n}_{\eta}(\varphi)\mathbin{\square}
     \operatorname{tr}^{\mathsf n}_{\eta}(\psi)
   \quad(\square\in\{\land,\lor,\to\}),\\
 \operatorname{tr}^{\mathsf n}_{\eta}(\forall v\,\varphi)
   &\defeq\forall x_{k(\eta)}\bigl(Q_{\operatorname{srt}(v)}(x_{k(\eta)})\to
     \operatorname{tr}^{\mathsf n}_{\eta\langle v\rangle}(\varphi)\bigr),\\
 \operatorname{tr}^{\mathsf n}_{\eta}(\exists v\,\varphi)
   &\defeq\exists x_{k(\eta)}\bigl(Q_{\operatorname{srt}(v)}(x_{k(\eta)})\land
     \operatorname{tr}^{\mathsf n}_{\eta\langle v\rangle}(\varphi)\bigr).
\end{aligned}
\]
Put \(\operatorname{tr}^{\mathsf n}\defeq\operatorname{tr}^{\mathsf n}_{\varnothing}\)
and \(\operatorname{tr}^{\mathsf n}[\Gamma]\defeq
\{\operatorname{tr}^{\mathsf n}(\gamma):\gamma\in\Gamma\}\).  Like
\(\operatorname{tr}_{\mathrm{MS}}\), this is a metalinguistic map, not an
object-language operator.
\index{guard translation!parameter form}
\end{hbdefinition}

Because \(x_{k(\eta)}\notin\operatorname{ran}\eta\), the function
\(\eta\langle v\rangle\) is again injective, so the recursion stays inside
environments.  Rebinding a variable already in \(\dom\eta\) discards its old
value: inside \(\forall v\,\forall v\,\varphi\) the inner occurrence of
\(v\) refers to the inner quantifier, as the many-sorted clauses require.

\begin{hblemma}{thm:p02-w02b-p02-033-sentences}{proved=annex:p02-w02b-p02-033-sentences}{Parameter names are fresh, and translations are sentences}
\begin{enumerate}
\item For distinct \(v,w\in\mathrm{Var}_{\mathrm{MS}}\),
      \(\mathsf n_v\ne\mathsf n_w\), and no \(\mathsf n_v\) lies in \(W\).
\item For every environment \(\eta\), every term \(t\) and every well-sorted
      formula \(\varphi\), the variables of \(t^{\eta}\) and the free
      variables of \(\operatorname{tr}^{\mathsf n}_{\eta}(\varphi)\) lie in
      \(\operatorname{ran}\eta\).  In particular
      \(\operatorname{tr}^{\mathsf n}(\varphi)\) is an
      \(\Lang_{\mathrm{MS}}^{\mathsf n}\)-sentence.
\end{enumerate}
\end{hblemma}

\subsection{The sort conventions of the parameter form}
\label{sec:p02-w02b-p02-033-conventions}

The parameter form uses the same sort theory as the variable form.  The theory
\(\mathsf{Ax}_{\mathrm{sort}}(\Lang_{\mathrm{MS}})\) consists of exactly the
four families of \cref{def:p02-w02-p02-038-axiom-families}: nonemptiness
\((\mathrm{Ne})\), one sentence per sort; constant typing \((\mathrm{Ct})\),
one per constant; function closure \((\mathrm{Cl})\), one per function symbol
of positive arity; and disjointness \((\mathrm{Dj})\), one per pair of
distinct sorts.  Each family is a set of first-order sentences, whatever its
cardinality, and the one-sorted consequence relation admits premise sets of
any size; so no step requires \(\mathsf{Sort}\), \(\mathcal F\) or
\(\mathcal R\) to be countable.  The premises contain no coverage axiom, and a
one-sorted model of the premises may have elements in no \(Q_{\mathsf s}\).
When \(\mathsf{Sort}=\{\mathsf s_1,\ldots,\mathsf s_k\}\) is finite, the
coverage variant adds the sentence \(\mathsf{Cov}(\Lang_{\mathrm{MS}})\) of
\cref{def:p02-w02-p02-038-coverage}.  When \(\mathsf{Sort}\) is infinite
there is no coverage variant, of any cardinality:
\cref{thm:p02-w02-p02-038-coverage-not-elementary} shows that no set of
\(\Lang_{\mathrm{MS}}^{\flat}\)-sentences expresses coverage over the models
of \(\mathsf{Ax}_{\mathrm{sort}}\), and its proof places no bound on
\(\mathsf{Sort}\), \(\mathcal F\) or \(\mathcal R\) (the compactness theorem
it invokes, \cref{thm:fol-compactness}, holds for every signature under the
Axiom of Choice, which that theorem assumes).

\begin{hbdefinition}{def:p02-w02b-p02-033-induced}{Induced sorted assignment}
Let \(\mathcal B\) be an \(\Lang_{\mathrm{MS}}^{\mathsf n}\)-structure, let
\(\eta\) be an environment and let \(\sigma\) be a one-sorted assignment into
\(\mathcal B\).  Then \(\sigma\) \emph{respects} \(\eta\) if
\(\sigma(\eta(v))\in Q_{\operatorname{srt}(v)}^{\mathcal B}\) for every
\(v\in\dom\eta\).  The \emph{induced sorted assignment} is the function
\(\theta_{\eta,\sigma}\) on \(\mathrm{Var}_{\mathrm{MS}}\) with
\[
 \theta_{\eta,\sigma}(v)\defeq
 \begin{cases}
  \sigma(\eta(v)) & v\in\dom\eta,\\
  \mathsf n_v^{\mathcal B} & v\notin\dom\eta.
 \end{cases}
\]
The restriction \(\mathcal B^{\downarrow}\) is the restriction of the
\(\Lang_{\mathrm{MS}}^{\flat}\)-reduct of \(\mathcal B\)
(\cref{def:p02-w02-p02-038-restriction}); it is defined when \(\mathcal B\)
satisfies \((\mathrm{Ct})\) and \((\mathrm{Cl})\).
\end{hbdefinition}

\begin{hbdefinition}{def:p02-w02b-p02-033-tagged}{Tagged parameter structures}
Let \(\mathcal M\) be a strict \(\Lang_{\mathrm{MS}}\)-structure and \(g\) a
total assignment into it.  Put
\(D_0\defeq\{(\mathsf s,a):\mathsf s\in\mathsf{Sort},\ a\in M_{\mathsf s}\}\)
and \(\star\defeq(\mathsf{Sort},\mathsf{Sort})\).  Both structures below
interpret \(Q_{\mathsf s}\) as \(\{\mathsf s\}\times M_{\mathsf s}\), a
constant \(c:\longrightarrow\mathsf s\) as \((\mathsf s,c^{\mathcal M})\),
\(\mathsf n_v\) as \((\operatorname{srt}(v),g(v))\), and a relation symbol
\(P:\mathsf s_1\times\cdots\times\mathsf s_n\) as
\(\{((\mathsf s_1,b_1),\ldots,(\mathsf s_n,b_n)):(b_1,\ldots,b_n)\in
P^{\mathcal M}\}\).  For
\(f:\mathsf s_1\times\cdots\times\mathsf s_n\longrightarrow\mathsf s\) with
\(n\geq1\), both send \(((\mathsf s_1,b_1),\ldots,(\mathsf s_n,b_n))\) with
every \(b_i\in M_{\mathsf s_i}\) to
\((\mathsf s,f^{\mathcal M}(b_1,\ldots,b_n))\).  They differ in the domain and
in the value at every other input:
\begin{enumerate}
\item \(\mathcal M^{\star}_{g}\) has domain \(D_0\cup\{\star\}\) and sends
      every other input to \(\star\);
\item \(\mathcal M^{\circ}_{g}\) has domain \(D_0\) and sends every other
      input of a symbol of codomain \(\mathsf s\) to
      \((\mathsf s,g(x^{\mathsf s}))\).
\end{enumerate}
The \(\Lang_{\mathrm{MS}}^{\flat}\)-reduct of \(\mathcal M^{\star}_{g}\) is
the structure \(\mathcal M^{\sharp}\), and that of \(\mathcal M^{\circ}_{g}\)
is \(\mathcal M^{\sharp}_{g}\), of
\cref{def:p02-w02-p02-038-transport}.
\index{many-sorted first-order logic!tagged structure}
\end{hbdefinition}

\subsection{The correspondence in both directions}
\label{sec:p02-w02b-p02-033-correspondence}

\begin{hblemma}{thm:p02-w02b-p02-033-isomorphism}{proved=annex:p02-w02b-p02-033-isomorphism}{Sorted isomorphisms preserve satisfaction}
Let \(\mathcal M\) and \(\mathcal N\) be sorted
\(\Lang_{\mathrm{MS}}\)-structures and let
\(\pi=(\pi_{\mathsf s})_{\mathsf s\in\mathsf{Sort}}\) be a family of
bijections \(\pi_{\mathsf s}:M_{\mathsf s}\to N_{\mathsf s}\) such that
\(\pi_{\mathsf s}(c^{\mathcal M})=c^{\mathcal N}\),
\(\pi_{\mathsf s}(f^{\mathcal M}(b_1,\ldots,b_n))=
f^{\mathcal N}(\pi_{\mathsf s_1}(b_1),\ldots,\pi_{\mathsf s_n}(b_n))\), and
\((b_1,\ldots,b_n)\in P^{\mathcal M}\) iff
\((\pi_{\mathsf s_1}(b_1),\ldots,\pi_{\mathsf s_n}(b_n))\in P^{\mathcal N}\),
for all declarations and arguments.  For a total assignment \(g\) into
\(\mathcal M\) let \(\pi g\) send \(v\) to
\(\pi_{\operatorname{srt}(v)}(g(v))\).  Then
\(\llbracket t\rrbracket^{\mathcal N}_{\pi g}=
\pi_{\mathsf s}(\llbracket t\rrbracket^{\mathcal M}_{g})\) for every term
\(t\) of sort \(\mathsf s\), and \(\mathcal M\entails\varphi[g]\) iff
\(\mathcal N\entails\varphi[\pi g]\) for every well-sorted formula
\(\varphi\).
\end{hblemma}

\begin{hbtheorem}{thm:p02-w02b-p02-033-restriction}{proved=annex:p02-w02b-p02-033-restriction}{Restriction correspondence for the parameter form}
Let \(\mathcal B\) be an \(\Lang_{\mathrm{MS}}^{\mathsf n}\)-structure
satisfying \((\mathrm{Ct})\), \((\mathrm{Cl})\) and
\(\mathsf{ParGuard}(\Lang_{\mathrm{MS}})\), let \(\eta\) be an environment, and
let \(\sigma\) be an assignment into \(\mathcal B\) that respects \(\eta\).
Write \(\theta\defeq\theta_{\eta,\sigma}\).
\begin{enumerate}
\item Every carrier \(Q_{\mathsf s}^{\mathcal B}\) of
      \(\mathcal B^{\downarrow}\) is nonempty; \(\mathcal B^{\downarrow}\) is
      strict if \(\mathcal B\entails(\mathrm{Dj})\); and \(\theta\) is a
      total assignment into \(\mathcal B^{\downarrow}\).
\item For every term \(t\) of sort \(\mathsf s\),
      \(\llbracket t^{\eta}\rrbracket^{\mathcal B}_{\sigma}=
      \llbracket t\rrbracket^{\mathcal B^{\downarrow}}_{\theta}\in
      Q_{\mathsf s}^{\mathcal B}\).
\item For every well-sorted formula \(\varphi\),
      \(\mathcal B,\sigma\entails\operatorname{tr}^{\mathsf n}_{\eta}(\varphi)\)
      iff \(\mathcal B^{\downarrow}\entails\varphi[\theta]\).
\end{enumerate}
\end{hbtheorem}

\begin{hbtheorem}{thm:p02-w02b-p02-033-tagging}{proved=annex:p02-w02b-p02-033-tagging}{Tagging correspondence for the parameter form}
Let \(\mathcal M\) be strict, let \(g\) be a total assignment into it, and let
\(\mathcal B\) be \(\mathcal M^{\star}_{g}\) or \(\mathcal M^{\circ}_{g}\).
\begin{enumerate}
\item \(\mathcal B\) satisfies \(\mathsf{Ax}_{\mathrm{sort}}(\Lang_{\mathrm{MS}})\)
      and \(\mathsf{ParGuard}(\Lang_{\mathrm{MS}})\).  The point \(\star\) of
      \(\mathcal M^{\star}_{g}\) lies in no \(Q_{\mathsf s}\).  If
      \(\mathsf{Sort}\) is finite, \(\mathcal M^{\circ}_{g}\) satisfies
      \(\mathsf{Cov}(\Lang_{\mathrm{MS}})\) and \(\mathcal M^{\star}_{g}\)
      does not.
\item Let \(\eta\) be an environment and \(\sigma\) an assignment into
      \(\mathcal B\) that respects \(\eta\).  Each
      \(\theta_{\eta,\sigma}(v)\) has the form
      \((\operatorname{srt}(v),a_v)\) with
      \(a_v\in M_{\operatorname{srt}(v)}\); put
      \(g_{\eta,\sigma}(v)\defeq a_v\).  Then for every well-sorted formula
      \(\varphi\),
      \(\mathcal B,\sigma\entails\operatorname{tr}^{\mathsf n}_{\eta}(\varphi)\)
      iff \(\mathcal M\entails\varphi[g_{\eta,\sigma}]\).
\item In particular, for every assignment \(\sigma\) into \(\mathcal B\),
      \(\mathcal B,\sigma\entails\operatorname{tr}^{\mathsf n}(\varphi)\) iff
      \(\mathcal M,g\entails\varphi\).
\end{enumerate}
\end{hbtheorem}

\begin{hbtheorem}{thm:p02-w02b-p02-033-reduction}{proved=annex:p02-w02b-p02-033-reduction}{Reduction of strict many-sorted consequence for every strict signature}
Let \(\Lang_{\mathrm{MS}}\) be a strict many-sorted signature, with
\(\mathsf{Sort}\), \(\mathcal F\) and \(\mathcal R\) of any cardinality, and
let \(\Gamma\cup\{\varphi\}\) be a set of well-sorted formulas.
\begin{enumerate}
\item \(\Gamma\entails_{\mathrm{MS}}\varphi\) iff
      \(\mathsf{Ax}_{\mathrm{sort}}(\Lang_{\mathrm{MS}})\cup
      \mathsf{ParGuard}(\Lang_{\mathrm{MS}})\cup
      \operatorname{tr}^{\mathsf n}[\Gamma]\entails_{\mathrm{FOL}}
      \operatorname{tr}^{\mathsf n}(\varphi)\).
\item If \(\mathsf{Sort}\) is finite, the same holds with
      \(\mathsf{Cov}(\Lang_{\mathrm{MS}})\) added to the premises.
\item Both equivalences remain true when the family \((\mathrm{Ne})\) is
      removed from the premises.
\end{enumerate}
The premises of clause~1 contain no coverage axiom.  No clause uses the Axiom
of Choice.
\index{many-sorted first-order logic!reduction}
\end{hbtheorem}

Clause~1 and \cref{thm:p02-w02b-p02-033-sentences} make many-sorted
consequence, with free variables, a consequence between one-sorted sentences.
The countermodel construction is the tagged structure with a junk point of
\cref{thm:p02-many-sorted-translation}, now carrying the values of the
assignment in its parameter names instead of in a one-sorted assignment.

\subsection{The variable form and where it stops}
\label{sec:p02-w02b-p02-033-variable-form}

\begin{hbdefinition}{def:p02-w02b-p02-033-renaming}{Variable renamings and the variable form}
A \emph{variable renaming} is any function
\(\iota:\mathrm{Var}_{\mathrm{MS}}\to\mathsf{Var}\).  For a term \(t\),
\(t^{\iota}\) replaces each variable \(v\) by \(\iota(v)\).  The map
\(\operatorname{tr}_{\iota}\) leaves \(\bottom\) unchanged, sends
\(P(t_1,\ldots,t_n)\) to \(P(t_1^{\iota},\ldots,t_n^{\iota})\) and \(t=u\) to
\(t^{\iota}=u^{\iota}\), commutes with \(\neg,\land,\lor,\to\), and satisfies
\[
\begin{aligned}
 \operatorname{tr}_{\iota}(\forall v\,\varphi)&\defeq
  \forall\iota(v)\bigl(Q_{\operatorname{srt}(v)}(\iota(v))\to
  \operatorname{tr}_{\iota}(\varphi)\bigr),\\
 \operatorname{tr}_{\iota}(\exists v\,\varphi)&\defeq
  \exists\iota(v)\bigl(Q_{\operatorname{srt}(v)}(\iota(v))\land
  \operatorname{tr}_{\iota}(\varphi)\bigr).
\end{aligned}
\]
Put \(\mathsf{VarGuard}_{\iota}\defeq
\{Q_{\operatorname{srt}(v)}(\iota(v)):v\in\mathrm{Var}_{\mathrm{MS}}\}\).  When
\(\iota\) is the fixed injection of \cref{def:p02-many-sorted-signature},
\(\operatorname{tr}_{\iota}=\operatorname{tr}_{\mathrm{MS}}\) and
\(\mathsf{VarGuard}_{\iota}=\mathsf{VarGuard}_{\mathrm{MS}}(\Lang_{\mathrm{MS}})\).
For an \(\Lang_{\mathrm{MS}}^{\flat}\)-structure \(\mathcal A\) and an
assignment \(\sigma\) into it, \(\mathcal A_{\sigma}^{\iota}\) is the
\(\Lang_{\mathrm{MS}}^{\mathsf n}\)-expansion of \(\mathcal A\) with
\(\mathsf n_v^{\mathcal A_{\sigma}^{\iota}}\defeq\sigma(\iota(v))\).
\end{hbdefinition}

\begin{hbtheorem}{thm:p02-w02b-p02-033-variable-form}{proved=annex:p02-w02b-p02-033-variable-form}{The variable form for every injective renaming}
Let \(\Lang_{\mathrm{MS}}\) be a strict signature with \(\mathcal F\) and
\(\mathcal R\) of any cardinality, and let \(\iota\) be an injective variable
renaming.
\begin{enumerate}
\item \emph{Bridge.}  For every \(\Lang_{\mathrm{MS}}^{\flat}\)-structure
      \(\mathcal A\), every assignment \(\sigma\) and every well-sorted
      \(\varphi\): \(\mathcal A,\sigma\entails\operatorname{tr}_{\iota}(\varphi)\)
      iff \(\mathcal A_{\sigma}^{\iota},\rho\entails
      \operatorname{tr}^{\mathsf n}(\varphi)\) for one, equivalently every,
      assignment \(\rho\); and \(\sigma\) satisfies
      \(\mathsf{VarGuard}_{\iota}\) iff
      \(\mathcal A_{\sigma}^{\iota}\entails\mathsf{ParGuard}(\Lang_{\mathrm{MS}})\).
\item \emph{Consequence.}  \(\Gamma\entails_{\mathrm{MS}}\varphi\) iff
      \(\mathsf{Ax}_{\mathrm{sort}}\cup\mathsf{VarGuard}_{\iota}\cup
      \operatorname{tr}_{\iota}[\Gamma]\entails_{\mathrm{FOL}}
      \operatorname{tr}_{\iota}(\varphi)\); if \(\mathsf{Sort}\) is finite,
      the same holds with \(\mathsf{Cov}\) added to the premises.
\item \emph{Restriction and tagging.}  If
      \(\mathcal A\entails\mathsf{Ax}_{\mathrm{sort}}\) and \(\sigma\)
      satisfies \(\mathsf{VarGuard}_{\iota}\), then \(\mathcal A^{\downarrow}\)
      is strict, \(\sigma\circ\iota\) is a total assignment into it,
      and \(\mathcal A,\sigma\entails\operatorname{tr}_{\iota}(\varphi)\) iff
      \(\mathcal A^{\downarrow}\entails\varphi[\sigma\circ\iota]\).
      If \(\mathcal M\) is strict, \(g\) is total, and \(g^{\sharp}\) sends
      \(\iota(v)\) to \((\operatorname{srt}(v),g(v))\) and every variable
      outside \(\operatorname{ran}\iota\) to \(\star\), then
      \(\mathcal M^{\sharp}\entails\mathsf{Ax}_{\mathrm{sort}}\),
      \(g^{\sharp}\) satisfies \(\mathsf{VarGuard}_{\iota}\), and
      \(\mathcal M,g\entails\varphi\) iff
      \(\mathcal M^{\sharp},g^{\sharp}\entails\operatorname{tr}_{\iota}(\varphi)\).
\end{enumerate}
Consequently the statements of \cref{thm:p02-many-sorted-translation},
\cref{thm:p02-w02-p02-038-strict-correspondence} and
\cref{thm:p02-w02-p02-038-coverage-reduction} hold for every strict signature
whose sorted variables are injected into \(\mathsf{Var}\), whatever the
cardinalities of \(\mathcal F\) and \(\mathcal R\).
\end{hbtheorem}

An injective renaming exists exactly when \(\mathrm{Var}_{\mathrm{MS}}\) is
countable.  If \(\mathsf{Sort}\) is at most countable, enumerated by an
injection \(e:\mathsf{Sort}\to\N\), and each stock is presented as
\(V_{\mathsf s}=\{x^{\mathsf s}_m:m\in\N\}\) without repetition, then
\(\iota(x^{\mathsf s}_m)\defeq x_{2^{e(\mathsf s)}(2m+1)-1}\) is one, defined
without choice.  With uncountably many sorts the variable form fails, for
every renaming.

\begin{hbtheorem}{thm:p02-w02b-p02-033-no-renaming}{proved=annex:p02-w02b-p02-033-no-renaming}{No exact variable form with uncountably many sorts}
\begin{enumerate}
\item If \(\mathsf{Sort}\) is uncountable, no variable renaming is
      injective.
\item If a variable renaming \(\iota\) is not injective, there is a
      well-sorted formula \(\varphi\), either \(\bottom\) or an equation
      between two variables, with
      \[
       \varnothing\not\entails_{\mathrm{MS}}\varphi
       \qquad\text{and}\qquad
       \mathsf{Ax}_{\mathrm{sort}}\cup\mathsf{VarGuard}_{\iota}
       \entails_{\mathrm{FOL}}\operatorname{tr}_{\iota}(\varphi).
      \]
\end{enumerate}
Hence, when \(\mathsf{Sort}\) is uncountable, the equivalence of clause~2 of
\cref{thm:p02-w02b-p02-033-variable-form} fails for every renaming of the
sorted variables into the countable one-sorted supply, and the reduction of
\cref{thm:p02-w02b-p02-033-reduction} is the one available.
\end{hbtheorem}

The failure is a failure of the variable supply, not of the sort axioms.  A
renaming that merges two variables of one sort makes their equation valid; a
renaming that merges variables of two sorts makes the guards contradict
\((\mathrm{Dj})\).  The parameter form avoids both, because free variables
become distinct constants, of which a signature may have any number, while
each formula, being finite, needs only finitely many one-sorted variables for
its bound occurrences.

\subsection{Example and limitations}
\label{sec:p02-w02b-p02-033-example}

\begin{hbexample}{def:p02-w02b-p02-033-example}{Uncountably many sorts}
Let \(\mathsf{Sort}=\mathbb R\), and for each real \(r\) declare a constant
\(e_r:\longrightarrow r\), with no other symbols.  Write \(y^r,z^r\) for two
distinct variables of sort \(r\) besides \(x^r\).  By
\cref{thm:p02-w02b-p02-033-no-renaming} the sorted variables admit no
injective renaming into \(\mathsf{Var}\), so only the parameter form applies.
For \(\psi\defeq\forall y^r\,\forall z^r\,(y^r=z^r)\) the recursion runs
through \(\varnothing\), then \(\{y^r\mapsto x_0\}\), then
\(\{y^r\mapsto x_0,z^r\mapsto x_1\}\):
\[
 \operatorname{tr}^{\mathsf n}(\psi)=
 \forall x_0\bigl(Q_r(x_0)\to\forall x_1(Q_r(x_1)\to x_0=x_1)\bigr).
\]
A rebinding resets the variable: with
\(\chi\defeq\forall y^r\,\exists y^r\,(y^r=e_r)\) the inner step starts from
\(\{y^r\mapsto x_0\}\), takes \(k=1\), and gives
\(\operatorname{tr}^{\mathsf n}(\chi)=\forall x_0\bigl(Q_r(x_0)\to
\exists x_1(Q_r(x_1)\land x_1=e_r)\bigr)\).  A free variable becomes its
name: \(\operatorname{tr}^{\mathsf n}(x^r=e_r)\) is
\(\mathsf n_{x^r}=e_r\).

Now \(\{\psi\}\entails_{\mathrm{MS}}x^r=e_r\), since \(\psi\) says that the
carrier of \(r\) has one element.  On the one-sorted side, a model of the
premises satisfies \(Q_r(\mathsf n_{x^r})\) by the parameter guard and
\(Q_r(e_r)\) by \((\mathrm{Ct})\); instantiating
\(\operatorname{tr}^{\mathsf n}(\psi)\) at these two elements gives
\(\mathsf n_{x^r}=e_r\), as \cref{thm:p02-w02b-p02-033-reduction} predicts.
Conversely \(\varnothing\not\entails_{\mathrm{MS}}x^r=e_r\): let
\(M_r=\{r\}\times\{0,1\}\) and \(e_r^{\mathcal M}=(r,0)\) for every real
\(r\), and \(g(v)=(\operatorname{srt}(v),1)\) for every sorted variable
\(v\).  The structure \(\mathcal M^{\star}_{g}\) has domain
\(\{(r,(r,i)):r\in\mathbb R,\ i\in\{0,1\}\}\cup\{\star\}\), satisfies all four
sort families (uncountably many \((\mathrm{Ne})\) and \((\mathrm{Dj})\)
sentences) and the parameter guard, and makes
\(\mathsf n_{x^r}=e_r\) false.  No coverage sentence exists here, as
\(\mathsf{Sort}\) is infinite.
\end{hbexample}

\begin{hbwarning}{def:p02-w02b-p02-033-limits}{What the parameter form does and does not give}
\begin{itemize}
\item The parameter form changes the shape of the translation: free sorted
      variables become constants of \(\Lang_{\mathrm{MS}}^{\mathsf n}\), and
      the assignment guard becomes a set of sentences.  The variable form of
      \cref{thm:p02-many-sorted-translation} survives only when the sorted
      variables inject into \(\mathsf{Var}\)
      (\cref{thm:p02-w02b-p02-033-variable-form,thm:p02-w02b-p02-033-no-renaming}).
\item The coverage variant needs finitely many sorts in either form
      (\cref{thm:p02-w02-p02-038-coverage-not-elementary}); the cardinality of
      \(\mathcal F\) and \(\mathcal R\) does not affect this.
\item Both forms are semantic.  A correspondence between a many-sorted proof
      calculus and a one-sorted one would need a translation of derivations,
      which is not given here.
\item The theorems concern the strict convention only.  The empty-sort and
      overlapping-sort reductions of
      \cref{sec:p02-w02-p02-038-empty,sec:p02-w02-p02-038-overlap} use the
      variable form and its injection.
\end{itemize}
\end{hbwarning}

\emph{Relationships.}  Whenever an injective renaming exists, clause~2 of
\cref{thm:p02-w02b-p02-033-variable-form} shows that the consequence clause
of \cref{thm:p02-many-sorted-translation} and clause~1 of
\cref{thm:p02-w02b-p02-033-reduction} assert the same thing.  The parameter
form reuses the restriction of \cref{def:p02-w02-p02-038-restriction} and the
tagged structures of \cref{def:p02-w02-p02-038-transport}.  The parameter
guard does for the parameter form what
\cref{thm:p02-w02-p02-038-guard-nonempty} records for the assignment guard:
it makes \((\mathrm{Ne})\) redundant.  Parameter names for elements are also
the device of the diagram method used with compactness
(\cref{thm:fol-compactness}); here they name the values of sorted variables
rather than elements of a given structure.

Signatures with uncountably many sorts occur naturally.  The internal
language of a category with finite products uses a many-sorted signature
whose sorts are the objects of the category and whose function symbols are
its morphisms; it is a set exactly when the category is small
\cite[Section~2.3, p.~10]{Pitts1995Draft}, so a small category with
uncountably many objects yields uncountably many sorts.  The same source
first fixes a disjoint countably infinite variable set for each sort, as
\cref{def:p02-w02b-p02-033-variables} does \cite[Section~2, p.~5]{Pitts1995Draft},
and then works with one fixed countably infinite variable set whose members
receive sorts only inside finite contexts
\cite[Section~2.1, Definition~2.1.1, p.~6]{Pitts1995Draft}.  Under that
second convention the variable supply does not grow with the number of sorts;
in the parameter form the finite environments \(\eta\) do the corresponding
work for bound variables, and the names \(\mathsf n_v\) carry the free ones.
